\documentclass[10pt,a4paper]{amsart}
\usepackage[margin=19mm]{geometry}
\usepackage{amsmath,amssymb,mathtools}
\usepackage[scr=rsfs,cal=euler]{mathalfa}
\usepackage{xfrac}
\usepackage[unicode,hidelinks,bookmarks=false]{hyperref}
\newcommand{\ie}{i.e.\ }
\newcommand{\cf}{cf.\ }
\newcommand{\resp}{respectively\ }
\newcommand{\Subsection}{\subsection}
\providecommand{\nobreakdash}{\nobreak\hbox{-}}
\setlength{\emergencystretch}{2em}
\multlinegap0pt
\usepackage{bm}
\usepackage{bbm}
\usepackage{dsfont}
\usepackage{xspace}
\usepackage{cancel}
\usepackage{mathtools}
\usepackage{amsthm}
\makeatletter
\@ifundefined{Theorem}{\newtheorem{Theorem}{Theorem}}{}
\makeatother
\let\iso=\cong
	\let\Union=\bigcup
	\let\Sect=\bigcap
	\let\Dirsum=\bigoplus
	\let\Tensor=\bigotimes
	%
	%------------------------------------------------
	%
	\let\to=\rightarrow
	\let\To=\longrightarrow
	\def\Implies{\ifmmode\Longrightarrow \else
		\unskip${}\Longrightarrow{}$\ignorespaces\fi}
\title[On the Linearity of Squarefree Powers of Edge Ideals]{On the Linearity of Squarefree Powers of Edge Ideals}
\author{Francesco Navarra, Ayesha Asloob Qureshi, Naoki Terai}
\date{Source: arXiv:2607.00838v1 (CC BY 4.0)}
\begin{document}
\maketitle

\noindent\emph{Source and licence.} On the Linearity of Squarefree Powers of Edge Ideals. Version arXiv:2607.00838v1 (CC BY 4.0), distributed under the Creative Commons Attribution 4.0 International licence, as recorded in the arXiv licence metadata for this exact version and shown on the abstract page https://arxiv.org/abs/2607.00838v1. Full rights notice: licenses/arxiv-2607.00838-CC-BY-4.0.txt.\par\smallskip

\noindent\textit{Editorial scope.} Counted unit: the Theorem whose source label is \texttt{Thm: Facet, pure, CM} in the article (source0.tex lines 914--935, source label \texttt{Thm: Facet, pure, CM}), delivered together with the article's own complete original proof of it (source0.tex lines 937--1036, one \texttt{proof} environment of 100 lines). No mathematics is written, repaired or re-derived: statement and proof are the article's own text line for line. Required context retained uncounted: none. Every definition, display and label of those lines is retained unchanged. Omitted from this package and not redistributed: 1--913, 936--936, 1037--2730. Nothing inside a retained excerpt is omitted or altered: every retained body is delivered exactly as the article's lines are written, so no omission receipt of any kind accompanies this package. Citations: the retained excerpts carry no citation command at all, so no bibliography entry of the article is transcribed and no reference is redistributed. Reference adaptation: none was needed and none was made; every reference inside the delivered unit resolves inside it, so no unresolved reference or citation is produced. Printed slips kept literal: no printed word, symbol, hypothesis or number of the counted statement or of its proof was altered. Layout and class: the article's own class is \texttt{amsart} with packages including amscd, amsmath, amsthm, amssymb, mathtools, geometry, epsfig, rawfonts, enumerate, graphics, multirow, xspace, graphicx, mathrsfs; the wrapper is the lane's pinned \texttt{amsart} editorial wrapper at 10pt on A4, which restates the theorem-like environments the retained text needs and adds the article's own macro definitions that the retained text needs (\texttt{\textbackslash iso}), with extra wrapper packages (bm, bbm, dsfont, xspace, cancel, mathtools). The delivered numbering is the unit's own. Assets: no retained span contains a figure environment, a graphics-inclusion command, a PDF-page inclusion, a table or a TikZ picture; no image file is copied into this package, the package declares no asset dependency and no image file is redistributed with it. Licence: version arXiv:2607.00838v1 of this article is distributed under the Creative Commons Attribution 4.0 International licence, as recorded in the arXiv licence metadata for this exact version and shown on the abstract page https://arxiv.org/abs/2607.00838v1. A full rights notice is delivered at licenses/arxiv-2607.00838-CC-BY-4.0.txt. Characters: every retained excerpt is pure ASCII, so the delivered engine, XeLaTeX with its default Latin Modern text font, reports no missing character.
\begin{Theorem}\label{Thm: Facet, pure, CM}
    Let $\Delta$ be a 1-dimensional flag simplicial complex on $[n]$ and $2\leq p\leq \nu(\overline{\Delta})$. Assume that $\dim S/I_{\Delta}^{[p]}=2p$. Then
    \begin{enumerate}
   \item The $(2p-1)$-dimensional facets of $\Delta^{[p]}$ are given by 
   \[
   \mathcal{F}_{2p-1}(\Delta^{[p]})= 
   \begin{cases} 
       \mathcal{F}( \Delta, K_{1,2p-1}) \cup \mathcal{F}( \Delta, K_{3,2p-3}) \cup \ldots \cup \mathcal{F}( \Delta, K_{p,p}) & \text{if } p \text{ is odd } \\
     \mathcal{F}( \Delta, K_{1,2p-1}) \cup \mathcal{F}( \Delta, K_{3,2p-3}) \cup \ldots \cup \mathcal{F}( \Delta, K_{p-1,p+1})   & \text{if } p \text{ is even}
   \end{cases}
\]
   
   \item Then $\Delta^{[p]}$ is pure if and only if one of the following holds
     \begin{enumerate}
       \item[(i)] $\Delta \iso K_{2k+1,2\ell+1}$ and $n=(2k+1)+(2\ell+1)\ge 2p$,
       \item[(ii)] $\Delta \iso K_{2k+1,2\ell}$ and $\ell \geq p$,
        \item[(iii)] $\Delta \iso K_{2k,2\ell}$ and $k,\ell \geq p$.
    \end{enumerate}

\item Then  $\Delta^{[p]}$ is Cohen-Macaulay if and only if $\Delta \iso K_{1,n-1}$ and $n>2p$.
    \end{enumerate}
\end{Theorem}

\begin{proof}

We divide the proof into three parts, corresponding to the statements above.

\textbf{(1)} First, we show by induction on $p$ that $\Delta_F$ is bipartite for any facet $F$ of $\Delta^{[p]}$ of dimension $2p-1$. \\
Let $F$ be a $(2p-1)-$dimensional facet of $\Delta^{[p]}$. %In the case $p=1$, there is nothing to prove and,
In the case $p=2$, $\Delta_F$ is trivially bipartite since $\Delta_F$ does not contain any triangle. Now assume $p \ge 3$. Suppose by contradiction that $\Delta_F$ is not bipartite. Then $\Delta_F$ contains an odd cycle $C_{2s+1}$ as an induced subgraph, for some $s\geq 1$.

\textit{Case 1.} Suppose that $2s+1 = 2p-1$.
Set $F \setminus V(C_{2s+1}) = \{v\}$.
Since $\Delta$ does not contain triangles, there exists $w \in V(C_{2s+1})$ such that $\{v,w\} \notin \Delta$.
One can see that
$$
\left( \prod_{u \in C_{2s+1} \setminus \{w\}} u \right)\in I_{\Delta}^{[p-1]}.
$$
Hence,
$$
vw \left( \prod_{u \in C_{2s+1} \setminus \{w\}} u \right)\in I_{\Delta}^{[p]}.
$$
Therefore, $F$ cannot be a facet of $\Delta^{[p]}$, which is a contradiction.

\textit{Case 2.} Suppose that $2s+1 \le 2p-3$.
Then there exist $v_1, v_2, v_3$ in $F \setminus V(C_{2s+1})$.
Since $\Delta$ does not contain triangles, we may assume that $\{v_1,v_2\} \notin \Delta$.
By induction on the number of vertices, we have
$$
\left(\prod_{u \in F \setminus \{v_1, v_2\}} u\right) \in I_{\Delta}^{[p-1]}.
$$
Hence
$$
v_1 v_2 \left(\prod_{u \in F \setminus \{v_1, v_2\}} u\right) \in I_{\Delta}^{[p]}.
$$
Therefore, $F$ cannot be a facet of $\Delta^{[p]}$, which is a contradiction.

Now we know that $\Delta_F$ is bipartite and is contained, as a spanning subgraph, in a complete bipartite graph $K_{\ell,m}$, where $\ell$ and $m$ are either both even or both odd.
If both $\ell$ and $m$ are even, the complement graph $\overline{K_{\ell,m}}$ has a $p$-matching. Hence so does $\overline{\Delta_F}$. Therefore, $F$ cannot be a facet of $\Delta^{[p]}$.
If both $\ell$ and $m$ are odd and $\Delta_F$ is not the complete bipartite graph $K_{\ell,m}$, then the complement graph $\overline{\Delta_F}$ has a $p$ a $p$-matching. Hence $F$ cannot be a facet of $\Delta^{[p]}$.
Thus, if $F$ is a facet of $\Delta^{[p]}$, then $\Delta_F$ must be the complete bipartite graph $K_{\ell,m}$ with $\ell$ and $m$ odd.

On the other hand, $\overline{K_{\ell,m}}$ with $\ell$ and $m$ odd does not have a $p$-matching. Hence, if $\Delta_F$ is $K_{\ell,m}$ with $\ell$ and $m$ odd, then $F$ is a facet of $\Delta^{[p]}$.

\textbf{(2)}
We assume that $\Delta^{[p]}$ is pure.
We show that $\Delta$ is bipartite.
Suppose that $\Delta$ contains an odd cycle $C_{2s+1}$ with $s \ge 2$.
Then we can take $W \subset V$ with $|W| = 3$ such that $\Delta_W$ is the disjoint union of
an edge and an isolated vertex. However, this cannot be contained as an induced subgraph in a complete bipartite graph $K_{\ell,m}$ with $\ell + m = 2p$ and $\ell, m$ odd. Hence we get a contradiction and, then, $\Delta$ has to be bipartite.

Next, we show that $\Delta$ is complete bipartite.
Suppose that $\Delta$ is contained in a complete bipartite graph $K_{s,t}$ with
$X \sqcup Y$ as vertex set partition, and that $\Delta$ is not complete bipartite.
We may assume that $|X| \ge 2$. Let $\{x, y\} \notin \Delta$, $\{x', y\} \in \Delta$, where $x, x' \in X$, $y \in Y$.
Take subsets $X' \subset X$, $Y' \subset Y$ such that $x, x' \in X'$, $y \in Y'$, and $|X'| + |Y'| = 2p-1$.
Then $\Delta_{X' \cup Y'}$ is neither a complete bipartite graph nor an edgeless graph (that is, a graph in which no vertices are connected by edges). Hence $\Delta_{X' \cup Y'}$ cannot be an induced subgraph of a complete bipartite graph.
Thus $X' \cup Y'$ is not contained in a facet of dimension $2p-1$.
Therefore, $\Delta$ is complete bipartite.

Since $p \leq \nu(\overline{\Delta})$, we have $n \ge 2p$.

If $\Delta \cong K_{2k+1,2\ell}$ with $\ell < p$ and $k > 0$, then 
$K_{2p-2\ell-1,2\ell}$ cannot be an induced subgraph of the form 
$K_{2s+1, 2t+1}$ with $(2s+1) + (2t+1) = 2p$.
Hence $V(K_{2p-2\ell-1,2\ell})$ is a facet of $\Delta^{[p]}$ of dimension $2p-2$. Therefore $\Delta^{[p]}$ is not pure of dimension $2p-1$.
Hence $\ell \ge p$ is necessary.

If $\Delta \cong K_{1,2\ell}$ with $\ell < p$, then $2\ell + 1 \le 2p-1$,
which contradicts the assumption that $p \leq \nu(\overline{\Delta})$.
Hence $\ell \ge p$ is necessary.

If $\Delta \cong K_{2k,2\ell}$ with $k \le \ell$ and $k < p$,
then $K_{2k,2p-2k-1}$ cannot be an induced subgraph of the form 
$K_{2s+1, 2t+1}$ with $(2s+1) + (2t+1) = 2p$.
Hence $V(K_{2k,2\ell})$ is a facet of $\Delta^{[p]}$ of dimension $2p-2$. Therefore $\Delta^{[p]}$ is not pure of dimension $2p-1$.
Hence $k, \ell \ge p$ is necessary.

Next, we show the converse.
For any facet $F$ of $\Delta^{[p]}$ with $|F| = 2p-1$,
$\Delta_F$ is a complete bipartite graph of the form $K_{2s,2t+1}$ with $2s + (2t+1) = 2p-1$, which is an induced subgraph of $K_{2s+1,2t+1}$,
whose vertex set is a facet of $\Delta^{[p]}$. Thus $\Delta^{[p]}$ is pure.

\textbf{(3)}
Assume that $\Delta^{[p]}$ is Cohen--Macaulay.
Then $\Delta^{[p]}$ is pure.
By (2), $\Delta$ is isomorphic to $K_{s,t}$ for some $s,t > 0$ with $s + t = n$.

Assume that $s,t > 1$ and that both $s$ and $t$ are odd.
Take $3 \le \ell \le s$ and $3 \le m \le t$ such that $\ell + m = 2p+2$.
Then there exist $F \in \Delta^{[p]}$ such that $\Delta_F \cong K_{\ell-2,m}$
and $G \in \Delta^{[p]}$ such that $\Delta_G \cong K_{\ell,m-2}$.
Then $F$ and $G$ cannot be connected in codimension one.

Assume that $s \ge 2p$ is even and $t > 1$.
Take $3 \le \ell \le s$ and $3 \le m \le t$ such that $\ell + m = 2p+2$.
Then there exist $F \in \Delta^{[p]}$ such that $\Delta_F \cong K_{2p-1,1}$
and $G \in \Delta^{[p]}$ such that $\Delta_G \cong K_{2p-3,3}$.
Then $F$ and $G$ cannot be connected in codimension one.

If $\Delta^{[p]}$ is Cohen--Macaulay, then $\Delta \cong K_{1,n-1}$.
If $n = 2p$, then $\nu(\overline{\Delta}) = p-1$, which contradicts the assumption that $\nu(\overline{\Delta}) \ge p$. Therefore $n > 2p$.
\end{proof}

\end{document}
